% ============================================================================
%  Informe técnico - Robot móvil diferencial con PIC16F877A y DRV8833
%  Taller Práctico II - Microcontroladores y Sistemas Embebidos
%  Universidad de Ibagué
%
%  Compilar:  latexmk -pdf main.tex      (o pdflatex dos veces)
%  En Overleaf: subir la carpeta completa y compilar main.tex con pdfLaTeX.
% ============================================================================
\documentclass[11pt,letterpaper]{article}

% ---------- Idioma y codificación ----------
\usepackage[utf8]{inputenc}
\usepackage[T1]{fontenc}
\usepackage[spanish,es-tabla,es-nodecimaldot]{babel}
\usepackage{lmodern}
\usepackage{microtype}

% ---------- Página ----------
\usepackage[margin=2.5cm]{geometry}
\usepackage{fancyhdr}
\usepackage{titlesec}
\usepackage{enumitem}
\setlist{itemsep=2pt, topsep=4pt}

% ---------- Matemáticas y unidades ----------
\usepackage{amsmath,amssymb}
\usepackage{siunitx}
\sisetup{output-decimal-marker={.}, group-separator={\,}, per-mode=symbol}

% ---------- Tablas y figuras ----------
\usepackage{booktabs}
\usepackage{tabularx}
\usepackage{array}
\usepackage{multirow}
\usepackage{float}
\usepackage{graphicx}
\usepackage{caption}
\captionsetup{font=small, labelfont=bf}
\usepackage{tikz}
\usetikzlibrary{arrows.meta, positioning, calc, fit, babel}  % babel: evita conflicto con el español

% ---------- Código ----------
\usepackage{xcolor}
\usepackage{listings}

% ---------- Enlaces (cargar al final) ----------
\usepackage[hidelinks]{hyperref}

% ============================================================================
%  Configuración general
% ============================================================================
% ============================================================================
%  config.tex - Datos del informe, colores, estilos y macros
% ============================================================================

% ---------- Datos del documento (editar aquí) ----------
\newcommand{\titulo}{Robot móvil diferencial con PIC16F877A y DRV8833}
\newcommand{\subtitulo}{Taller Práctico II: Control dual de motores y robótica móvil diferencial}
\newcommand{\asignatura}{Microcontroladores y Sistemas Embebidos}
\newcommand{\programa}{Programa de Ingeniería Electrónica}
\newcommand{\facultad}{Facultad de Ingeniería}
\newcommand{\universidad}{Universidad de Ibagué}
\newcommand{\docente}{\pendiente{Nombre del docente}}
\newcommand{\fecha}{Octubre de 2026}
\newcommand{\versiondoc}{0.1 (borrador)}

% ---------- Colores ----------
\definecolor{acento}{RGB}{0,61,121}       % azul institucional aprox.
\definecolor{acentoclaro}{RGB}{220,232,245}
\definecolor{gris}{RGB}{110,110,110}
\definecolor{codigofondo}{RGB}{247,247,247}
\definecolor{codigocomentario}{RGB}{0,128,0}
\definecolor{codigoclave}{RGB}{0,0,170}
\definecolor{codigocadena}{RGB}{163,21,21}
\definecolor{pendientecolor}{RGB}{200,0,0}

% ---------- Marcas de contenido pendiente ----------
% Uso: \pendiente{qué falta}. Buscar "\pendiente" para ver lo que falta.
\newcommand{\pendiente}[1]{\textcolor{pendientecolor}{\textit{[#1]}}}

% Recuadro para fotografías o capturas que aún no se han agregado.
% Uso: \figurapendiente{descripción}{etiqueta}
\newcommand{\figurapendiente}[2]{%
  \begin{figure}[H]
    \centering
    \fbox{\parbox[c][4cm][c]{0.8\textwidth}{\centering\pendiente{Agregar: #1}}}
    \caption{#1}
    \label{#2}
  \end{figure}}

% ---------- Encabezado y pie ----------
\pagestyle{fancy}
\fancyhf{}
\fancyhead[L]{\small\color{gris}\asignatura}
\fancyhead[R]{\small\color{gris}\universidad}
\fancyfoot[L]{\small\color{gris}Versión \versiondoc}
\fancyfoot[R]{\small\thepage}
\renewcommand{\headrulewidth}{0.4pt}
\setlength{\headheight}{14pt}

% ---------- Títulos de sección ----------
\titleformat{\section}{\Large\bfseries\color{acento}}{\thesection.}{0.6em}{}
\titleformat{\subsection}{\large\bfseries\color{acento}}{\thesubsection}{0.6em}{}
\titlespacing*{\section}{0pt}{18pt}{8pt}
\titlespacing*{\subsection}{0pt}{12pt}{6pt}

% ---------- Tablas ----------
\renewcommand{\arraystretch}{1.25}
\newcolumntype{L}{>{\raggedright\arraybackslash}X}
\newcolumntype{C}{>{\centering\arraybackslash}X}

% ---------- Listados de código ----------
\lstdefinestyle{c-xc8}{
  language=C,
  basicstyle=\ttfamily\footnotesize,
  backgroundcolor=\color{codigofondo},
  commentstyle=\color{codigocomentario},
  keywordstyle=\color{codigoclave}\bfseries,
  stringstyle=\color{codigocadena},
  morekeywords={__interrupt, volatile, __delay_ms},
  numbers=left,
  numberstyle=\tiny\color{gris},
  numbersep=8pt,
  frame=single,
  rulecolor=\color{gris!40},
  breaklines=true,
  breakatwhitespace=true,
  showstringspaces=false,
  tabsize=4,
  xleftmargin=1.5em,
  framexleftmargin=1.5em,
  captionpos=t
}
\lstset{style=c-xc8}
\renewcommand{\lstlistingname}{Código}
\renewcommand{\lstlistlistingname}{Índice de códigos}

% ---------- Metadatos del PDF ----------
\hypersetup{
  pdftitle={\titulo},
  pdfsubject={\subtitulo},
  pdfkeywords={PIC16F877A, DRV8833, PWM, robótica móvil, XC8}
}


% ============================================================================
\begin{document}

% ============================================================================
%  Portada
% ============================================================================
\begin{titlepage}
  \centering
  % Logo de la universidad (opcional): guardar como figuras/logo-unibague.png
  % y descomentar la siguiente línea.
  % \includegraphics[width=0.35\textwidth]{figuras/logo-unibague.png}\par\vspace{1cm}

  {\Large\scshape \universidad\par}
  \vspace{0.2cm}
  {\large \facultad\ --- \programa\par}
  \vspace{0.2cm}
  {\large \asignatura\par}

  \vspace{2.5cm}
  {\color{acento}\rule{\textwidth}{1.2pt}\par}
  \vspace{0.5cm}
  {\huge\bfseries\color{acento} \titulo\par}
  \vspace{0.4cm}
  {\Large \subtitulo\par}
  \vspace{0.5cm}
  {\color{acento}\rule{\textwidth}{1.2pt}\par}

  \vspace{2cm}
  {\large\textbf{Informe técnico}\par}

  \vspace{1.5cm}
  \begin{tabular}{rl}
    \textbf{Integrantes:} & Brahiam Alberto Campiño Romero \\
                          & \pendiente{Integrante 2} \\
                          & \pendiente{Integrante 3} \\[6pt]
    \textbf{Docente:}     & \docente \\
  \end{tabular}

  \vfill
  {\large Ibagué, Tolima --- \fecha\par}
  \vspace{0.3cm}
  {\small\color{gris} Versión \versiondoc\par}
\end{titlepage}


\tableofcontents
\clearpage

% ============================================================================
%  Resumen
% ============================================================================
\section*{Resumen}
\addcontentsline{toc}{section}{Resumen}

Se diseñó e implementó el control de una plataforma móvil diferencial de dos
motorreductores DC con un PIC16F877A y el doble puente H DRV8833. La velocidad
de cada rueda se controla por PWM de hardware (CCP1 y CCP2 sincronizados con
Timer2) a \textbf{\SI{1.25}{\kilo\hertz} con 1000 niveles de ciclo de trabajo
($\approx$10~bits)}, y el sentido de giro y el frenado dinámico se gobiernan
con dos líneas digitales de dirección.

El firmware, escrito en C con MPLAB XC8, está organizado en tres capas
(hardware, control por motor y navegación) y ofrece dos modos seleccionables en
compilación: la secuencia automática de validación del taller (avance, reversa,
giro y frenado con rampas) y el control remoto desde un celular por Bluetooth
(HC-05/HC-06, UART a 9600~baudios, con interrupción de recepción y control de
velocidad en 11 niveles).

Durante el diseño se identificó que la frecuencia de \SI{1.2}{\kilo\hertz}
pedida no es alcanzable con cristal de \SI{20}{\mega\hertz} y prescaler 1:16
(exige $PR2 = 259 > 255$); se justifica el uso de \SI{1.25}{\kilo\hertz}.
También se resolvieron la ausencia del pulso de habilitación nSLEEP y del giro
en el código base, y la diferencia entre el cristal de la lista de materiales
(\SI{4}{\mega\hertz}) y el del código (\SI{20}{\mega\hertz}).

\medskip
\noindent\textbf{Palabras clave:} PIC16F877A, DRV8833, PWM, Timer2, robótica
móvil diferencial, puente H, UART, Bluetooth, MPLAB XC8, sistemas embebidos.

% ============================================================================
%  1. Introducción y objetivos
% ============================================================================
\section{Introducción y objetivos}
\label{sec:introduccion}

Un robot móvil diferencial se desplaza y orienta variando solo la velocidad y
el sentido de sus dos ruedas motrices: misma velocidad y sentido produce
traslación, sentidos opuestos producen rotación sobre su eje. Esta
arquitectura es la base de robots de servicio, AGV industriales y plataformas
de competencia, y exige controlar dos motores DC de forma independiente y
sincronizada.

\subsection*{Objetivo general}
Implementar el control de velocidad, sentido de giro y frenado de una
plataforma móvil diferencial con un PIC16F877A y un driver DRV8833, validado en
simulación y en hardware real.

\subsection*{Objetivos específicos}
\begin{enumerate}
  \item Configurar y sincronizar los módulos CCP1 (RC2) y CCP2 (RC1) en modo
        PWM con Timer2 como base de tiempo única.
  \item Integrar el DRV8833 para gobernar velocidad, sentido de giro y frenado
        dinámico de dos motores DC.
  \item Aplicar normas de montaje en protoboard: masa común, habilitación por
        nSLEEP y supresión de ruido inductivo ($dI/dt$).
  \item Desarrollar firmware modular en MPLAB XC8 para la navegación de la
        plataforma.
  \item \emph{(Extensión del grupo)} Agregar control remoto desde un celular
        mediante un módulo Bluetooth serie.
\end{enumerate}

% ============================================================================
%  2. Arquitectura del sistema
% ============================================================================
\section{Arquitectura del sistema}
\label{sec:arquitectura}

El sistema separa la lógica (\SI{5}{\volt} regulados por el LM7805) de la
potencia (batería de hasta \SI{10}{\volt} hacia el VM del DRV8833), unidas solo
por una masa común. El PIC genera cinco señales hacia el driver: dos PWM de
hardware para la velocidad, dos líneas de dirección y la habilitación nSLEEP
(figura~\ref{fig:arquitectura}).

\begin{figure}[H]
  \centering
  % ============================================================================
%  Diagrama de bloques del sistema (TikZ)
%  Se incluye desde secciones/03-arquitectura.tex con % ============================================================================
%  Diagrama de bloques del sistema (TikZ)
%  Se incluye desde secciones/03-arquitectura.tex con \input{figuras/arquitectura}
% ============================================================================
\begin{tikzpicture}[
    font=\small,
    >={Stealth[length=2.2mm]},
    bloque/.style={draw=gris, rounded corners=3pt, line width=0.6pt,
                   minimum width=3.2cm, minimum height=1.1cm, align=center},
    principal/.style={draw=acento, fill=acentoclaro, rounded corners=3pt,
                      line width=1.2pt, align=left, inner sep=6pt},
    driver/.style={draw=gris, rounded corners=3pt, line width=0.6pt,
                   align=left, inner sep=6pt},
    flecha/.style={->, draw=gris, line width=0.6pt},
    etiqueta/.style={font=\footnotesize, text=gris}
  ]

  % ---------- Microcontrolador ----------
  \node[principal, minimum width=5.2cm, minimum height=9.6cm, anchor=north]
       (pic) at (0,0) {};
  \node[font=\bfseries\large, anchor=north] at ($(pic.north)+(0,-0.25)$) {PIC16F877A};
  \node[anchor=north west, align=left, font=\footnotesize]
       at ($(pic.north west)+(0.2,-1.0)$) {%
       RC2 / CCP1: PWM izquierdo\\[2pt]
       RD0: dirección izquierda\\[2pt]
       RC1 / CCP2: PWM derecho\\[2pt]
       RD1: dirección derecha\\[2pt]
       RD2: nSLEEP};

  % ---------- Columna izquierda ----------
  \node[bloque, anchor=north] (bat) at (-5.6,-0.1)
       {\textbf{Batería}\\ \footnotesize\textcolor{gris}{máx.\ \SI{10}{\volt} DC}};
  \node[bloque, below=0.75cm of bat] (reg)
       {\textbf{LM7805}\\ \footnotesize\textcolor{gris}{\SI{5}{\volt} para la lógica}};
  \node[bloque, below=0.75cm of reg] (xtal)
       {\textbf{Cristal}\\ \footnotesize\textcolor{gris}{\SI{20}{\mega\hertz} (o \SI{4}{\mega\hertz})}};
  \node[bloque, below=0.75cm of xtal] (pk)
       {\textbf{PICkit 3}\\ \footnotesize\textcolor{gris}{programación ICSP}};
  \node[bloque, below=0.75cm of pk] (bt)
       {\textbf{HC-05 / HC-06}\\ \footnotesize\textcolor{gris}{Bluetooth, 9600 baud}};

  % ---------- Columna derecha ----------
  \node[driver, minimum width=4.2cm, minimum height=4.2cm, anchor=north]
       (drv) at (6.4,-0.1) {};
  \node[font=\bfseries\large, anchor=north] at ($(drv.north)+(0,-0.2)$) {DRV8833};
  \node[anchor=north west, align=left, font=\footnotesize]
       at ($(drv.north west)+(0.15,-0.9)$) {%
       AIN1 / AIN2: rueda izq.\\[2pt]
       BIN1 / BIN2: rueda der.\\[2pt]
       nSLEEP: habilitación\\[2pt]
       \textcolor{gris}{VM: 2.7 a \SI{10.8}{\volt}}\\[2pt]
       \textcolor{gris}{\SI{1.5}{\ampere} por canal}};

  \node[bloque, minimum width=4.2cm, below=0.9cm of drv] (mizq)
       {\textbf{Motor izquierdo}\\ \footnotesize\textcolor{gris}{AOUT1 / AOUT2}};
  \node[bloque, minimum width=4.2cm, below=0.75cm of mizq] (mder)
       {\textbf{Motor derecho}\\ \footnotesize\textcolor{gris}{BOUT1 / BOUT2}};

  % ---------- Conexiones ----------
  % Potencia de motores
  \draw[flecha] (bat.north) -- ++(0,0.45) -| node[etiqueta, pos=0.25, above]
       {Vmotor $\leq$ \SI{10}{\volt}} (drv.north);
  \draw[flecha] (bat.south) -- (reg.north);

  % Lógica y periféricos hacia el PIC
  \draw[flecha] (reg.east)  -- node[etiqueta, above] {\SI{5}{\volt}} (reg.east -| pic.west);
  \draw[flecha] (xtal.east) -- (xtal.east -| pic.west);
  \draw[flecha] (pk.east)   -- (pk.east -| pic.west);
  \draw[<->, draw=gris, line width=0.6pt]
       (bt.east) -- node[etiqueta, above] {UART} (bt.east -| pic.west);

  % Pines del PIC alineados con cada flecha que llega
  \node[anchor=west, font=\footnotesize, text=gris] at (xtal.east -| pic.west)
       [xshift=0.2cm, yshift=0.15cm] {OSC1 / OSC2: cristal};
  \node[anchor=west, font=\footnotesize, text=gris] at (xtal.east -| pic.west)
       [xshift=0.2cm, yshift=-0.25cm] {Timer2: base de tiempo PWM};
  \node[anchor=west, font=\footnotesize, text=gris] at (pk.east -| pic.west)
       [xshift=0.2cm] {MCLR, RB6, RB7: ICSP};
  \node[anchor=west, font=\footnotesize, text=gris] at (bt.east -| pic.west)
       [xshift=0.2cm] {RC6 TX / RC7 RX: UART};

  % Señales de control hacia el driver (línea horizontal)
  \coordinate (sal) at ($(drv.west)+(0,0.6)$);
  \draw[flecha] (sal -| pic.east) -- node[etiqueta, above] {5 señales} (sal);

  % Salidas a los motores
  \draw[flecha] (drv.south) -- (mizq.north);
  \draw[flecha] ($(drv.east)+(0,-1.0)$) -- ++(0.4,0) |- (mder.east);
\end{tikzpicture}

% ============================================================================
\begin{tikzpicture}[
    font=\small,
    >={Stealth[length=2.2mm]},
    bloque/.style={draw=gris, rounded corners=3pt, line width=0.6pt,
                   minimum width=3.2cm, minimum height=1.1cm, align=center},
    principal/.style={draw=acento, fill=acentoclaro, rounded corners=3pt,
                      line width=1.2pt, align=left, inner sep=6pt},
    driver/.style={draw=gris, rounded corners=3pt, line width=0.6pt,
                   align=left, inner sep=6pt},
    flecha/.style={->, draw=gris, line width=0.6pt},
    etiqueta/.style={font=\footnotesize, text=gris}
  ]

  % ---------- Microcontrolador ----------
  \node[principal, minimum width=5.2cm, minimum height=9.6cm, anchor=north]
       (pic) at (0,0) {};
  \node[font=\bfseries\large, anchor=north] at ($(pic.north)+(0,-0.25)$) {PIC16F877A};
  \node[anchor=north west, align=left, font=\footnotesize]
       at ($(pic.north west)+(0.2,-1.0)$) {%
       RC2 / CCP1: PWM izquierdo\\[2pt]
       RD0: dirección izquierda\\[2pt]
       RC1 / CCP2: PWM derecho\\[2pt]
       RD1: dirección derecha\\[2pt]
       RD2: nSLEEP};

  % ---------- Columna izquierda ----------
  \node[bloque, anchor=north] (bat) at (-5.6,-0.1)
       {\textbf{Batería}\\ \footnotesize\textcolor{gris}{máx.\ \SI{10}{\volt} DC}};
  \node[bloque, below=0.75cm of bat] (reg)
       {\textbf{LM7805}\\ \footnotesize\textcolor{gris}{\SI{5}{\volt} para la lógica}};
  \node[bloque, below=0.75cm of reg] (xtal)
       {\textbf{Cristal}\\ \footnotesize\textcolor{gris}{\SI{20}{\mega\hertz} (o \SI{4}{\mega\hertz})}};
  \node[bloque, below=0.75cm of xtal] (pk)
       {\textbf{PICkit 3}\\ \footnotesize\textcolor{gris}{programación ICSP}};
  \node[bloque, below=0.75cm of pk] (bt)
       {\textbf{HC-05 / HC-06}\\ \footnotesize\textcolor{gris}{Bluetooth, 9600 baud}};

  % ---------- Columna derecha ----------
  \node[driver, minimum width=4.2cm, minimum height=4.2cm, anchor=north]
       (drv) at (6.4,-0.1) {};
  \node[font=\bfseries\large, anchor=north] at ($(drv.north)+(0,-0.2)$) {DRV8833};
  \node[anchor=north west, align=left, font=\footnotesize]
       at ($(drv.north west)+(0.15,-0.9)$) {%
       AIN1 / AIN2: rueda izq.\\[2pt]
       BIN1 / BIN2: rueda der.\\[2pt]
       nSLEEP: habilitación\\[2pt]
       \textcolor{gris}{VM: 2.7 a \SI{10.8}{\volt}}\\[2pt]
       \textcolor{gris}{\SI{1.5}{\ampere} por canal}};

  \node[bloque, minimum width=4.2cm, below=0.9cm of drv] (mizq)
       {\textbf{Motor izquierdo}\\ \footnotesize\textcolor{gris}{AOUT1 / AOUT2}};
  \node[bloque, minimum width=4.2cm, below=0.75cm of mizq] (mder)
       {\textbf{Motor derecho}\\ \footnotesize\textcolor{gris}{BOUT1 / BOUT2}};

  % ---------- Conexiones ----------
  % Potencia de motores
  \draw[flecha] (bat.north) -- ++(0,0.45) -| node[etiqueta, pos=0.25, above]
       {Vmotor $\leq$ \SI{10}{\volt}} (drv.north);
  \draw[flecha] (bat.south) -- (reg.north);

  % Lógica y periféricos hacia el PIC
  \draw[flecha] (reg.east)  -- node[etiqueta, above] {\SI{5}{\volt}} (reg.east -| pic.west);
  \draw[flecha] (xtal.east) -- (xtal.east -| pic.west);
  \draw[flecha] (pk.east)   -- (pk.east -| pic.west);
  \draw[<->, draw=gris, line width=0.6pt]
       (bt.east) -- node[etiqueta, above] {UART} (bt.east -| pic.west);

  % Pines del PIC alineados con cada flecha que llega
  \node[anchor=west, font=\footnotesize, text=gris] at (xtal.east -| pic.west)
       [xshift=0.2cm, yshift=0.15cm] {OSC1 / OSC2: cristal};
  \node[anchor=west, font=\footnotesize, text=gris] at (xtal.east -| pic.west)
       [xshift=0.2cm, yshift=-0.25cm] {Timer2: base de tiempo PWM};
  \node[anchor=west, font=\footnotesize, text=gris] at (pk.east -| pic.west)
       [xshift=0.2cm] {MCLR, RB6, RB7: ICSP};
  \node[anchor=west, font=\footnotesize, text=gris] at (bt.east -| pic.west)
       [xshift=0.2cm] {RC6 TX / RC7 RX: UART};

  % Señales de control hacia el driver (línea horizontal)
  \coordinate (sal) at ($(drv.west)+(0,0.6)$);
  \draw[flecha] (sal -| pic.east) -- node[etiqueta, above] {5 señales} (sal);

  % Salidas a los motores
  \draw[flecha] (drv.south) -- (mizq.north);
  \draw[flecha] ($(drv.east)+(0,-1.0)$) -- ++(0.4,0) |- (mder.east);
\end{tikzpicture}

  \caption{Arquitectura del sistema: control, potencia, actuadores y
           comunicaciones.}
  \label{fig:arquitectura}
\end{figure}

\subsection{Asignación de pines}

\begin{table}[H]
  \centering
  \caption{Asignación de pines del PIC16F877A y justificación.}
  \label{tab:pines}
  \small
  \begin{tabularx}{\textwidth}{@{}l l l L@{}}
    \toprule
    \textbf{Pin PIC} & \textbf{Función} & \textbf{Conectado a} & \textbf{Justificación} \\
    \midrule
    RC2 / CCP1 (17) & PWM motor izq.   & AIN1 & Único pin de salida del módulo CCP1; PWM por hardware sin carga de CPU. \\
    RC1 / CCP2 (16) & PWM motor der.   & BIN1 & Único pin de salida de CCP2; comparte Timer2 con CCP1, así ambas ruedas tienen la misma frecuencia y fase. \\
    RD0 (19)        & Dirección izq.   & AIN2 & Puerto D libre de funciones especiales que interfieran. \\
    RD1 (20)        & Dirección der.   & BIN2 & Ídem. \\
    RD2 (21)        & Habilitación     & nSLEEP & Permite dormir el driver hasta terminar la configuración y aplicar el pulso de activación. \\
    RC6 / TX (25)   & UART TX          & RXD HC-05 & Pin fijo del USART; un divisor 1k/2k baja \SI{5}{\volt} a $\approx$\SI{3.3}{\volt}. \\
    RC7 / RX (26)   & UART RX          & TXD HC-05 & Pin fijo del USART con interrupción de recepción. \\
    RB6, RB7, MCLR  & ICSP             & PICkit 3 & Reservados para programar en circuito sin retirar el PIC. \\
    \bottomrule
  \end{tabularx}
\end{table}

\subsection{Componentes principales}

PIC16F877A, módulo DRV8833, dos motorreductores DC (caja amarilla tipo TT),
módulo Bluetooth HC-05/HC-06, regulador LM7805, cristal de cuarzo, batería de
máximo \SI{10}{\volt} DC y programador PICkit~3.

% ============================================================================
%  3. Etapa de potencia: DRV8833
% ============================================================================
\section{Etapa de potencia: DRV8833}
\label{sec:potencia}

El DRV8833 integra dos puentes H de MOSFET que operan entre \SI{2.7}{\volt} y
\SI{10.8}{\volt} con hasta \SI{1.5}{\ampere} continuos por canal. Cada motor
se controla con dos entradas ($x$IN1, $x$IN2), según la tabla~\ref{tab:drv}.

\begin{table}[H]
  \centering
  \caption{Tabla de verdad del DRV8833 por canal.}
  \label{tab:drv}
  \begin{tabular}{@{}c c l l@{}}
    \toprule
    \textbf{$x$IN1} & \textbf{$x$IN2} & \textbf{Salidas ($x$OUT1 / $x$OUT2)} & \textbf{Modo} \\
    \midrule
    0 & 0 & Alta impedancia & Parada libre (\emph{coast}) \\
    1 & 0 & Vmotor / GND    & Avance \\
    0 & 1 & GND / Vmotor    & Reversa \\
    1 & 1 & GND / GND       & Freno dinámico (\emph{slow decay}) \\
    \bottomrule
  \end{tabular}
\end{table}

\subsection{Por qué DRV8833 y no L298N}

El L298N usa transistores bipolares en Darlington y pierde cerca de
\SI{2}{\volt} o más entre sus salidas; con una batería de 6 a \SI{10}{\volt}
eso es una fracción grande del voltaje del motor y se disipa como calor, por lo
que necesita disipador. El DRV8833 usa MOSFET de baja resistencia de
encendido: la caída es mucho menor, el motor recibe casi todo el voltaje de la
batería y el driver no requiere disipador volumétrico.

\subsection{Aplicación del PWM en cada sentido}

\begin{itemize}
  \item \textbf{Avance:} $x$IN2 = 0 fijo y PWM en $x$IN1. El driver alterna
        entre 10 (avance) y 00 (parada libre, decaimiento rápido). La velocidad
        es proporcional al ciclo de trabajo $D$.
  \item \textbf{Reversa:} $x$IN2 = 1 fijo y PWM en $x$IN1 con ciclo invertido
        ($1000 - \text{velocidad}$). El driver alterna entre 01 (reversa) y 11
        (freno, decaimiento lento). El tiempo en reversa es la fracción en bajo
        del PWM; por eso se invierte el ciclo: velocidad 700 $\rightarrow$
        duty 300 $\rightarrow$ 70\,\% del periodo en reversa.
  \item \textbf{Freno dinámico:} $x$IN1 = $x$IN2 = 1. Ambos terminales del
        motor quedan a GND, la fuerza contraelectromotriz circula por los
        MOSFET inferiores y el motor se detiene mucho más rápido que en parada
        libre.
\end{itemize}

\subsection{Habilitación nSLEEP}

Con nSLEEP en 0 el driver entra en bajo consumo y sus salidas quedan en alta
impedancia. El firmware lo mantiene en 0 mientras configura PWM y pines, y
luego aplica el flanco $0 \rightarrow 1$ con una espera de unos milisegundos
(tiempo de arranque del driver, del orden de \SI{1}{\milli\second}). Así,
ningún estado transitorio de los pines durante el reinicio del PIC mueve los
motores.

% ============================================================================
%  4. Montaje físico y normas de ensamble
% ============================================================================
\section{Montaje físico y normas de ensamble}
\label{sec:montaje}

Cada regla de montaje responde a un modo de falla concreto de los sistemas que
mezclan lógica y motores (tabla~\ref{tab:montaje}).

\begin{table}[H]
  \centering
  \caption{Normas de montaje físico y falla que previene cada una.}
  \label{tab:montaje}
  \small
  \begin{tabularx}{\textwidth}{@{}>{\raggedright\arraybackslash}p{2.4cm} L L@{}}
    \toprule
    \textbf{Regla} & \textbf{Implementación} & \textbf{Falla que evita} \\
    \midrule
    Masa común & GND de \SI{5}{\volt} y de Vmotor unidos en un solo punto cerca del DRV8833. &
      Sin referencia común el driver no interpreta los niveles lógicos del PIC; unirlas en varios puntos crea lazos de tierra por donde circula corriente de motor. \\
    Desacoplo del PIC & Cerámico de \SI{0.1}{\micro\farad} entre VDD y VSS (pines 11/12 y 32/31). &
      Picos de consumo del núcleo que causan reinicios por caída de VDD. \\
    Bus de potencia & Electrolítico de \SI{100}{\micro\farad} en paralelo con cerámico de \SI{0.1}{\micro\farad} entre VM y GND del driver. &
      El electrolítico entrega los picos de arranque e inversión; el cerámico filtra el ruido de alta frecuencia del PWM. \\
    Supresión en motores & Cerámico de \SI{0.1}{\micro\farad} entre terminales de cada motor. &
      Ruido de las escobillas y transitorios $dI/dt$ que se acoplan a la lógica. \\
    Regulación & LM7805 desde la batería para VDD = \SI{5}{\volt}. &
      La batería cae al arrancar los motores; el regulador mantiene estable la lógica. \\
    ICSP & MCLR, RB6 y RB7 libres hacia el conector del PICkit~3. &
      Permite reprogramar sin retirar el PIC del montaje. \\
    Reset & \SI{10}{\kilo\ohm} de MCLR a VDD y pulsador NA a GND. &
      MCLR flotante provoca reinicios aleatorios. \\
    Cableado & UTP calibre 24~AWG cortado a medida. &
      Cables largos actúan como antenas y elevan la inductancia de las rutas de corriente. \\
    \bottomrule
  \end{tabularx}
\end{table}

\subsection{Oscilador}

El cristal va lo más cerca posible de OSC1/OSC2, con dos capacitores a GND. La
hoja de datos del PIC16F877A recomienda 15 a \SI{33}{\pico\farad} para
cristales de 4 y \SI{20}{\mega\hertz}; el esquemático usa \SI{22}{\pico\farad}.
Los \SI{44}{\pico\farad} de la lista de materiales quedan fuera de ese rango y
pueden dificultar el arranque del oscilador.

\subsection{Nivel lógico del HC-05}

Su entrada RXD trabaja a \SI{3.3}{\volt}, por lo que la salida TX del PIC
(\SI{5}{\volt}) pasa por un divisor de \SI{1}{\kilo\ohm} y \SI{2}{\kilo\ohm}:
\begin{equation}
  V_{RXD} = \SI{5}{\volt}\cdot\frac{2}{1+2} \approx \SI{3.3}{\volt}
\end{equation}
En sentido contrario, el pin RX del PIC (entrada Schmitt trigger) requiere en
hoja de datos $0.8\,V_{DD} = \SI{4}{\volt}$ para leer un 1. Los
\SI{3.3}{\volt} del módulo suelen funcionar en la práctica; si aparecen errores
de recepción se agrega un convertidor de nivel.

\figurapendiente{Fotografía del montaje en protoboard}{fig:montaje}

% ============================================================================
%  5. Justificación matemática
% ============================================================================
\section{Justificación matemática}
\label{sec:calculos}

Con cristal de \SI{20}{\mega\hertz}, prescaler 1:16 y $PR2 = 249$, el PWM
opera a \textbf{\SI{1250}{\hertz} con 1000 niveles de ciclo de trabajo
(\num{9.97}~bits)}. Los \SI{1.2}{\kilo\hertz} pedidos no son alcanzables con
esa configuración.

\subsection{Ciclo de instrucción}

El PIC16F877A ejecuta una instrucción cada 4 ciclos de reloj:
\begin{align}
  T_{osc} &= \frac{1}{F_{osc}} = \frac{1}{\SI{20}{\mega\hertz}} = \SI{50}{\nano\second} \\
  T_{cy}  &= 4\,T_{osc} = \SI{200}{\nano\second}
  \qquad\Rightarrow\qquad F_{cy} = \SI{5}{\mega\hertz}
\end{align}

\subsection{Registro PR2 para 1.2 kHz}

El periodo del PWM lo fija Timer2, que se reinicia al igualar a PR2:
\begin{equation}
  T_{PWM} = (PR2 + 1)\cdot 4\,T_{osc}\cdot \text{Pre}
  \qquad\Rightarrow\qquad
  PR2 = \frac{F_{osc}}{4\cdot \text{Pre}\cdot F_{PWM}} - 1
  \label{eq:pr2}
\end{equation}
\begin{equation}
  PR2 = \frac{\num{20000000}}{4\cdot 16\cdot 1200} - 1 = 260.42 - 1 = 259.42
\end{equation}

PR2 es un registro de 8~bits (máximo 255), por lo que \textbf{259.42 no se
puede cargar}. La frecuencia mínima posible con \SI{20}{\mega\hertz} y
prescaler 1:16 es:
\begin{equation}
  F_{PWM,\min} = \frac{\num{20000000}}{4\cdot 16\cdot 256} = \SI{1220.7}{\hertz}
\end{equation}

El prescaler del Timer2 no puede superar 1:16 en el PIC16F877A, así que
\SI{1.2}{\kilo\hertz} exacto es inalcanzable con este cristal. Se elige
$\mathbf{PR2 = 249}$:
\begin{equation}
  F_{PWM} = \frac{\num{20000000}}{4\cdot 16\cdot 250} = \SI{1250}{\hertz}
  \qquad
  T_{PWM} = \SI{800}{\micro\second}
\end{equation}

El error frente a \SI{1.2}{\kilo\hertz} es de 4.2\,\%. Con $PR2 = 255$ se
obtendría \SI{1220.7}{\hertz} (1.7\,\% de error), pero se prefiere 249 porque
deja exactamente 1000 cuentas por periodo: el ciclo de trabajo se escribe
directamente en escala 0--1000 (décimas de porcentaje) sin conversiones.
Además, \SI{1.25}{\kilo\hertz} está en el rango típico de PWM para
motorreductores pequeños: la inductancia del motor filtra la corriente y el
control es estable a baja velocidad.

\subsection{Ciclo de trabajo y resolución de 10 bits}

El ciclo de trabajo es un valor de 10~bits: 8~bits en CCPR$x$L y 2~bits en
CCP$x$CON$\langle5{:}4\rangle$ (CCP$x$X, CCP$x$Y).
\begin{equation}
  T_{on} = DC\cdot T_{osc}\cdot \text{Pre} = DC\cdot \SI{50}{\nano\second}\cdot 16
         = DC\cdot \SI{0.8}{\micro\second}
\end{equation}
\begin{equation}
  \text{Pasos} = 4\,(PR2+1) = 1000
  \qquad
  \text{Resolución} = \log_2(1000) = \num{9.97}~\text{bits}
  \qquad
  \Delta D = \SI{0.1}{\percent}
\end{equation}

\begin{table}[H]
  \centering
  \caption{Distribución del ciclo de trabajo en los registros del módulo CCP.}
  \label{tab:duty}
  \begin{tabular}{@{}r r c r r@{}}
    \toprule
    \textbf{$DC$} & \textbf{CCPR$x$L} & \textbf{CCP$x$X:CCP$x$Y} & \textbf{$T_{on}$} & \textbf{Ciclo} \\
    \midrule
    0    & 0   & 00 & \SI{0}{\micro\second}   & 0\,\%   \\
    250  & 62  & 10 & \SI{200}{\micro\second} & 25\,\%  \\
    500  & 125 & 00 & \SI{400}{\micro\second} & 50\,\%  \\
    750  & 187 & 10 & \SI{600}{\micro\second} & 75\,\%  \\
    1000 & 250 & 00 & \SI{800}{\micro\second} & 100\,\% \\
    \bottomrule
  \end{tabular}
\end{table}

\subsection{Caso con cristal de 4 MHz}

Si se monta el cristal de \SI{4}{\mega\hertz} de la lista de materiales con la
configuración anterior, el PWM cae a
$\SI{4}{\mega\hertz} / (64 \cdot 250) = \SI{250}{\hertz}$: frecuencia audible y
par irregular a baja velocidad. Para conservar $PR2 = 249$ y la escala 0--1000
se usa prescaler 1:4:
\begin{equation}
  F_{PWM} = \frac{\num{4000000}}{4\cdot 4\cdot 250} = \SI{1000}{\hertz}
  \qquad
  T_{cy} = \SI{1}{\micro\second}
\end{equation}

\subsection{Velocidad de la UART (Bluetooth)}

Con BRGH = 1 (alta velocidad):
\begin{equation}
  \text{Baud} = \frac{F_{osc}}{16\,(SPBRG + 1)}
\end{equation}

\begin{table}[H]
  \centering
  \caption{Configuración de SPBRG para 9600 baudios.}
  \label{tab:uart}
  \begin{tabular}{@{}l r r r@{}}
    \toprule
    \textbf{Cristal} & \textbf{SPBRG} & \textbf{Baud real} & \textbf{Error vs 9600} \\
    \midrule
    \SI{20}{\mega\hertz} & 129 & 9615 & +0.16\,\% \\
    \SI{4}{\mega\hertz}  & 25  & 9615 & +0.16\,\% \\
    \bottomrule
  \end{tabular}
\end{table}

Un error menor a $\pm$2\,\% es aceptable para comunicación asíncrona de 8~bits.

% ============================================================================
%  6. Diseño del firmware
% ============================================================================
\section{Diseño del firmware}
\label{sec:firmware}

El firmware se organiza en tres capas: cada capa solo usa las funciones de la
capa inferior, de modo que cambiar un pin o el driver obliga a tocar solo la
capa de hardware. Dos directivas al inicio del archivo seleccionan en
compilación el cristal y el modo de operación, sin duplicar código
(código~\ref{lst:config}).

\begin{lstlisting}[caption={Selección de cristal y modo de operación.}, label={lst:config}, numbers=none]
#define CRISTAL_20MHZ   1   // 1 = 20 MHz (HS, 1:16) | 0 = 4 MHz (XT, 1:4)
#define MODO_BLUETOOTH  0   // 0 = secuencia del taller | 1 = control por celular
\end{lstlisting}

\subsection{Funciones por capa}

\begin{table}[H]
  \centering
  \caption{Funciones del firmware organizadas por capa.}
  \label{tab:funciones}
  \small
  \begin{tabularx}{\textwidth}{@{}l >{\ttfamily\raggedright\arraybackslash}p{4.6cm} L@{}}
    \toprule
    \textbf{Capa} & \textbf{\normalfont Función} & \textbf{Responsabilidad} \\
    \midrule
    Hardware     & PWM\_Init()                       & Configura TRIS, $PR2 = 249$, CCP1/CCP2 en PWM y Timer2; deja nSLEEP en 0. \\
    Hardware     & DRV8833\_Wake()                   & Pulso $0 \rightarrow 1$ en nSLEEP y espera de arranque del driver. \\
    Hardware     & Set\_Duty\_Motor\_Izq/\allowbreak Der()       & Satura a 1000 y reparte el valor de 10~bits entre CCPR$x$L y CCP$x$X:CCP$x$Y. \\
    Motor        & Motor\_Izquierdo/\allowbreak Derecho(speed, dir) & Fija $x$IN2 y aplica PWM directo (avance) o invertido (reversa). \\
    Motor        & Freno\_General()                  & $x$IN1 = $x$IN2 = 1 en ambos canales. \\
    Motor        & Parada\_Libre()                   & $x$IN1 = $x$IN2 = 0; estado seguro al arrancar. \\
    Navegación   & Mover(mov, v)                     & Traduce ADELANTE, ATRAS, IZQUIERDA, DERECHA o PARAR a órdenes por rueda. \\
    Navegación   & Rampa\_Subir/Bajar(mov)           & Acelera o desacelera de 0 a 1000 en pasos de 50 cada \SI{100}{\milli\second}. \\
    Comunicación & UART\_Init(), ISR                 & USART a 9600~baudios y recepción de comandos por interrupción. \\
    \bottomrule
  \end{tabularx}
\end{table}

\subsection{Cinemática diferencial}

\begin{table}[H]
  \centering
  \caption{Combinación de sentidos por rueda para cada movimiento.}
  \label{tab:cinematica}
  \begin{tabular}{@{}l l l l@{}}
    \toprule
    \textbf{Movimiento} & \textbf{Rueda izquierda} & \textbf{Rueda derecha} & \textbf{Resultado} \\
    \midrule
    Adelante  & Avance  & Avance  & Traslación hacia adelante \\
    Atrás     & Reversa & Reversa & Traslación hacia atrás \\
    Izquierda & Reversa & Avance  & Rotación antihoraria sobre su eje \\
    Derecha   & Avance  & Reversa & Rotación horaria sobre su eje \\
    Parar     & Freno   & Freno   & Detención con freno dinámico \\
    \bottomrule
  \end{tabular}
\end{table}

Se usa giro sobre el propio eje (ruedas en sentidos opuestos) en lugar de
frenar una sola rueda: el radio de giro es cero, la maniobra es simétrica y la
plataforma puede reorientarse en espacios reducidos.

\subsection{Modo 0: secuencia de validación del taller}

\begin{enumerate}
  \item Avance con rampa de aceleración (0 $\rightarrow$ 1000 en 21 pasos de
        \SI{100}{\milli\second}), \SI{1}{\second} a velocidad máxima y rampa de
        desaceleración.
  \item Freno dinámico durante \SI{500}{\milli\second}.
  \item Reversa con las mismas rampas y freno de \SI{500}{\milli\second}.
  \item Giro a la izquierda a velocidad 600 durante \SI{1}{\second} y freno de
        \SI{500}{\milli\second}.
  \item Giro a la derecha a velocidad 600 durante \SI{1}{\second}.
  \item Freno total durante \SI{2}{\second} y repetición del ciclo.
\end{enumerate}

Las rampas limitan la corriente de arranque de los motores: un escalón de 0 a
100\,\% exige al driver y a la batería la corriente de rotor bloqueado, y la
caída de voltaje resultante es la causa típica de reinicios del
microcontrolador, que la rúbrica penaliza.

\subsection{Modo 1: control remoto por Bluetooth}

\begin{table}[H]
  \centering
  \caption{Comandos aceptados por el modo Bluetooth.}
  \label{tab:comandos}
  \begin{tabular}{@{}l l@{}}
    \toprule
    \textbf{Comando} & \textbf{Acción} \\
    \midrule
    \texttt{F} / \texttt{B}     & Adelante / atrás \\
    \texttt{L} / \texttt{R}     & Giro izquierda / derecha \\
    \texttt{S}                  & Freno dinámico \\
    \texttt{0} \dots\ \texttt{9} & Velocidad 0\,\% \dots\ 90\,\% \\
    \texttt{q}                  & Velocidad 100\,\% \\
    \bottomrule
  \end{tabular}
\end{table}

Decisiones de diseño del modo remoto:
\begin{itemize}
  \item \textbf{Recepción por interrupción.} La ISR solo lee RCREG y actualiza
        dos variables \texttt{volatile} (comando y velocidad); el lazo
        principal aplica los cambios. Así no se pierden caracteres aunque el
        programa esté en un retardo, y la ISR es corta.
  \item \textbf{Recuperación de \emph{overrun}.} Si el búfer de dos bytes del
        USART se desborda (OERR), se reinicia CREN; sin esto, la recepción
        queda bloqueada para siempre.
  \item \textbf{Lectura atómica.} La velocidad es de 16~bits y el PIC la lee en
        dos accesos; el lazo deshabilita GIE durante la copia para no leer un
        valor a medio actualizar.
  \item \textbf{Freno antes de invertir.} Al cambiar de movimiento se aplica
        freno dinámico \SI{50}{\milli\second} antes del nuevo sentido, lo que
        evita el pico de corriente de invertir un motor en marcha.
  \item \textbf{Eco.} Cada cambio de movimiento devuelve su letra al celular
        como confirmación.
  \item \textbf{Compatibilidad.} Las letras coinciden con la aplicación
        \emph{Bluetooth RC Controller} (Android), que envía la letra mientras se
        mantiene la flecha y \texttt{S} al soltarla.
\end{itemize}

El código completo se presenta en el anexo~\ref{anx:codigo}.

% ============================================================================
%  7. Simulación en Proteus
% ============================================================================
\section{Simulación en Proteus}
\label{sec:simulacion}

Proteus no incluye un modelo del DRV8833 en su librería estándar, por lo que se
simula con un \textbf{L293D}, que tiene la misma tabla de verdad por canal
(IN1/IN2) y permite validar toda la lógica del firmware. Las diferencias de
potencia (caída de voltaje, corriente máxima) solo se evalúan en el montaje
real.

\begin{table}[H]
  \centering
  \caption{Elementos y configuración de la simulación.}
  \label{tab:proteus}
  \small
  \begin{tabularx}{\textwidth}{@{}l L L@{}}
    \toprule
    \textbf{Elemento} & \textbf{Conexión} & \textbf{Configuración} \\
    \midrule
    PIC16F877A       & Program File = \texttt{.hex} de MPLAB X & Processor Clock Frequency = \SI{20}{\mega\hertz} (o \SI{4}{\mega\hertz}) \\
    L293D            & IN1 = RC2, IN2 = RD0, IN3 = RC1, IN4 = RD1 & EN1 y EN2 a +\SI{5}{\volt}; VSS a +\SI{5}{\volt}; VS a la fuente de motores \\
    2 $\times$ MOTOR (DC) & OUT1/OUT2 y OUT3/OUT4 & Motor izquierdo y derecho \\
    OSCILLOSCOPE     & Canal A en RC2, canal B en RC1 & Medición de periodo y ciclo de trabajo \\
    VIRTUAL TERMINAL & TXD $\rightarrow$ RC7, RXD $\rightarrow$ RC6 & 9600~baudios, 8N1 (reemplaza al HC-05) \\
    RES + BUTTON     & MCLR con \SI{10}{\kilo\ohm} a +\SI{5}{\volt} & Reset \\
    \bottomrule
  \end{tabularx}
\end{table}

\subsection{Procedimiento de prueba}

\begin{enumerate}
  \item Compilar en MPLAB X con \texttt{MODO\_BLUETOOTH 0} y cargar el
        \texttt{.hex}.
  \item Verificar con el osciloscopio un periodo de \SI{800}{\micro\second} en
        RC2 y RC1, y el crecimiento del ciclo de trabajo durante las rampas.
  \item Observar el sentido de giro de los motores en cada etapa de la
        secuencia.
  \item Compilar con \texttt{MODO\_BLUETOOTH 1}, escribir \texttt{F},
        \texttt{B}, \texttt{L}, \texttt{R}, \texttt{S} y \texttt{0}--\texttt{9}
        en la terminal virtual y comprobar la respuesta y el eco.
\end{enumerate}

\figurapendiente{Esquemático del circuito en Proteus}{fig:proteus-esquematico}
\figurapendiente{Osciloscopio con PWM al 25\,\%, 50\,\% y 75\,\%}{fig:proteus-pwm}
\figurapendiente{Terminal virtual con comandos y eco}{fig:proteus-terminal}

% ============================================================================
%  8. Pruebas y resultados
% ============================================================================
\section{Pruebas y resultados}
\label{sec:pruebas}

Las tablas recogen el valor esperado de cada prueba según la
sección~\ref{sec:calculos} y dejan las columnas de medición para completar en
simulación y en hardware real.
% Para llenar: reemplazar cada \pendiente{...} por el valor medido y
% "Pendiente" por "Cumple" o "No cumple".

\subsection{Mediciones de PWM (cristal de 20 MHz)}

\begin{table}[H]
  \centering
  \caption{Mediciones eléctricas esperadas y obtenidas.}
  \label{tab:mediciones}
  \small
  \begin{tabularx}{\textwidth}{@{}L >{\raggedright\arraybackslash}p{3.0cm} c c c@{}}
    \toprule
    \textbf{Prueba} & \textbf{Esperado} & \textbf{Proteus} & \textbf{Hardware} & \textbf{Resultado} \\
    \midrule
    Frecuencia PWM en RC2               & \SI{1250}{\hertz}                 & \pendiente{} & \pendiente{} & Pendiente \\
    Frecuencia PWM en RC1               & \SI{1250}{\hertz}                 & \pendiente{} & \pendiente{} & Pendiente \\
    $T_{on}$ con $DC = 250$             & \SI{200}{\micro\second} (25\,\%) & \pendiente{} & \pendiente{} & Pendiente \\
    $T_{on}$ con $DC = 500$             & \SI{400}{\micro\second} (50\,\%) & \pendiente{} & \pendiente{} & Pendiente \\
    $T_{on}$ con $DC = 750$             & \SI{600}{\micro\second} (75\,\%) & \pendiente{} & \pendiente{} & Pendiente \\
    Duración de rampa 0 $\rightarrow$ 1000 & \SI{2.1}{\second} (21 pasos)  & \pendiente{} & \pendiente{} & Pendiente \\
    Voltaje en VM con motores girando   & $\leq$ \SI{10}{\volt}            & ---          & \pendiente{} & Pendiente \\
    VDD del PIC al arrancar motores     & \SI{5}{\volt} $\pm$ \SI{0.25}{\volt} & ---       & \pendiente{} & Pendiente \\
    Pulso en nSLEEP al encender         & Flanco $0 \rightarrow 1$ tras configurar PWM & \pendiente{} & \pendiente{} & Pendiente \\
    UART: comando \texttt{F} con eco    & Avance y eco \texttt{F}          & \pendiente{} & \pendiente{} & Pendiente \\
    \bottomrule
  \end{tabularx}
\end{table}

\subsection{Demostración funcional (rúbrica, criterio 4)}

\begin{table}[H]
  \centering
  \caption{Validación de maniobras en hardware real.}
  \label{tab:demostracion}
  \small
  \begin{tabularx}{\textwidth}{@{}l L L c@{}}
    \toprule
    \textbf{Maniobra} & \textbf{Comportamiento esperado} & \textbf{Observación} & \textbf{Resultado} \\
    \midrule
    Avance         & Ambas ruedas hacia adelante, aceleración suave & \pendiente{} & Pendiente \\
    Reversa        & Ambas ruedas hacia atrás, aceleración suave    & \pendiente{} & Pendiente \\
    Giro izquierda & Rotación antihoraria sobre su eje              & \pendiente{} & Pendiente \\
    Giro derecha   & Rotación horaria sobre su eje                  & \pendiente{} & Pendiente \\
    Freno          & Detención rápida sin reinicio del PIC          & \pendiente{} & Pendiente \\
    \bottomrule
  \end{tabularx}
\end{table}

% ============================================================================
%  9. Decisiones de diseño y discrepancias resueltas
% ============================================================================
\section{Decisiones de diseño y discrepancias resueltas}
\label{sec:decisiones}

Al contrastar la guía del taller con las hojas de datos se encontraron cinco
puntos que, sin resolverse, habrían afectado el funcionamiento o la evaluación
(tabla~\ref{tab:decisiones}).

\begin{table}[H]
  \centering
  \caption{Hallazgos y decisiones de diseño.}
  \label{tab:decisiones}
  \small
  \begin{tabularx}{\textwidth}{@{}c L L L@{}}
    \toprule
    \textbf{\#} & \textbf{Hallazgo} & \textbf{Impacto} & \textbf{Decisión} \\
    \midrule
    1 & \SI{1.2}{\kilo\hertz} con \SI{20}{\mega\hertz} y prescaler 1:16 exige $PR2 = 259.4$, fuera del rango de 8~bits. &
        Imposible de configurar. &
        $PR2 = 249 \rightarrow \SI{1250}{\hertz}$ (error 4.2\,\%), con escala exacta de 1000 pasos (sección~\ref{sec:calculos}). \\
    2 & La lista de materiales indica cristal de \SI{4}{\mega\hertz}; el código usa \SI{20}{\mega\hertz} (FOSC = HS). &
        Con \SI{4}{\mega\hertz} el PWM baja a \SI{250}{\hertz} y los retardos se alargan 5 veces. &
        Directiva \texttt{CRISTAL\_20MHZ} que ajusta FOSC, prescaler y SPBRG para ambos casos. \\
    3 & La rúbrica evalúa el pulso de activación en nSLEEP, pero el código base no maneja ese pin. &
        El driver depende de un pull-up del módulo o queda deshabilitado. &
        nSLEEP en RD2, en 0 durante la configuración y flanco $0 \rightarrow 1$ con \texttt{DRV8833\_Wake()}. \\
    4 & La rúbrica evalúa giro, pero la secuencia base solo tiene avance, reversa y freno. &
        Criterio de demostración incompleto. &
        Giro a la izquierda y a la derecha sobre su eje agregados a la secuencia. \\
    5 & Capacitores de cristal de \SI{44}{\pico\farad} en la lista de materiales. &
        Fuera del rango recomendado de 15 a \SI{33}{\pico\farad}; arranque del oscilador poco confiable. &
        Usar \SI{22}{\pico\farad}, como en el esquemático. \\
    \bottomrule
  \end{tabularx}
\end{table}

También se definió el sentido de giro según la convención física (izquierda =
rotación antihoraria vista desde arriba), de modo que el comando y el
movimiento observado coincidan durante la demostración.

% ============================================================================
%  10. Conclusiones y trabajo futuro
% ============================================================================
\section{Conclusiones y trabajo futuro}
\label{sec:conclusiones}

\subsection{Conclusiones}

\begin{itemize}
  \item El uso de CCP1 y CCP2 con Timer2 como base de tiempo única genera dos
        PWM de hardware de igual frecuencia y fase sin consumir tiempo de CPU,
        lo que deja el lazo principal libre para la navegación y la
        comunicación.
  \item La justificación matemática demostró que la especificación de
        \SI{1.2}{\kilo\hertz} no es alcanzable con \SI{20}{\mega\hertz} y
        prescaler 1:16; elegir $PR2 = 249$ da \SI{1.25}{\kilo\hertz} y una
        escala de 1000 pasos (0.1\,\%) que simplifica el firmware.
  \item Aplicar PWM directo en avance e invertido en reversa con el DRV8833
        permite controlar velocidad y sentido con solo dos líneas por motor, y
        el freno dinámico detiene la plataforma más rápido que la parada libre.
  \item Las rampas de aceleración, el freno previo a cada inversión y el
        desacoplo del bus de potencia atacan la causa principal de reinicios
        del microcontrolador: las caídas de voltaje por corriente de arranque.
  \item La arquitectura en capas y la selección de modo en compilación
        permitieron agregar el control Bluetooth sin modificar las funciones de
        control de motores.
\end{itemize}

\subsection{Competencias aplicadas}

Configuración de periféricos de microcontroladores (CCP, Timer2, USART,
interrupciones), electrónica de potencia para motores DC, diseño de firmware
embebido modular en C, comunicación serie inalámbrica, buenas prácticas de
integridad de señal y alimentación, simulación de circuitos en Proteus y
validación en hardware.

\subsection{Trabajo futuro}

\begin{itemize}
  \item Lazo cerrado de velocidad con encoders en las ruedas y control PI, para
        que el robot avance recto aunque los dos motores no sean idénticos.
  \item Seguridad por pérdida de enlace: detener el robot si no llegan
        comandos Bluetooth durante un tiempo definido.
  \item Giros suaves con velocidades distintas por rueda (giro con radio),
        además del giro sobre el eje.
  \item Medición del voltaje de batería con el ADC para alertar carga baja.
  \item Migración del montaje de protoboard a una PCB a partir del
        esquemático.
\end{itemize}

% ============================================================================
%  Referencias
% ============================================================================
\begin{thebibliography}{9}
\addcontentsline{toc}{section}{Referencias}

\bibitem{pic16f877a}
Microchip Technology Inc., \emph{PIC16F87XA Data Sheet: 28/40/44-Pin Enhanced
Flash Microcontrollers}, DS39582. Secciones de CCP/PWM, Timer2, USART y
oscilador. \url{https://www.microchip.com/en-us/product/PIC16F877A}

\bibitem{drv8833}
Texas Instruments, \emph{DRV8833 Dual H-Bridge Motor Driver}, hoja de datos.
\url{https://www.ti.com/product/DRV8833}

\bibitem{xc8}
Microchip Technology Inc., \emph{MPLAB XC8 C Compiler User's Guide for PIC
MCU}.

\bibitem{guia}
Universidad de Ibagué, \emph{Taller Práctico II: Control dual de motores y
robótica móvil diferencial}, Microcontroladores y Sistemas Embebidos, 2026.

\bibitem{hc05}
\emph{HC-05 / HC-06 Bluetooth Serial Module}, hoja de datos del fabricante del
módulo.

\end{thebibliography}


\appendix
% ============================================================================
%  Anexo A. Código fuente
%  El código se lee directamente de ../firmware/taller2_robot_diferencial.c,
%  así que cualquier cambio en el firmware aparece solo en el informe.
% ============================================================================
\section{Código fuente completo}
\label{anx:codigo}

El firmware está en el archivo \texttt{firmware/taller2\_robot\_diferencial.c}
del repositorio. Se compila con MPLAB XC8 para el PIC16F877A.

\lstinputlisting[caption={Firmware del robot móvil diferencial
(\texttt{taller2\_robot\_diferencial.c}).}, label={lst:firmware}]
{../firmware/taller2_robot_diferencial.c}


\end{document}
